\documentclass[eng]{tex/class/phynum-class}

\title{Chaos-based symmetric cryptography : a naive implementation for image encryption}
\shorttitle{Chaos-based symmetric cryptography}
       
\author[1]{DE SOUSA Maxime}
\affil[1]{Université Grenoble-Alpes, Physique parcours recherche fondamentale, Grenoble, France}

\firstauthor{DE SOUSA Maxime}

\abstract{With the recents developments of data transmission, encryption become a real society concern. Traditional encryption algorithm like the Advanced Encryption Standard (AES), aren't suitable to image encryption due to large computational time, or because of the redundancy of the information in the data. This lead us to think that it might exist better algorithm for image encryption. In this paper, we describe a naive implementation for image encryption using chaotic maps. Basic analysis methods, such as visual, intensity and correlation, are used to evaluate the proposed algorithm in regards of image encryption. We finally show that the proposed algorithm, can be a good basis to image encryption.}

\keywords{cryptography, chaos, dynamical systems, symmetric-key, encryption}

\begin{document}

%include title, authors, abstract, etc.
\maketitle
\thispagestyle{fancy}

%main
\section{Introduction}
	\firstword{T}{he} needs for encryption quickly emerged in human's history, mostly because of military operations or to exchange sensitive informations. However, nowadays, with the development of Internet, everybody needs to exchange encrypted data : texts, videos, images.. 
Most of the traditional, and accepted, symmetric algorithms (AES-128, RC6..) to encrypt data are block-cipher algorithm. This type of algorithm lead to large computational time to encrypt large images. Moreover, the redundancy of information in image data lead some encryption algorithm to show some weaknesses when it comes to image encryption.  

Chaos-based cryptography can be suitable to encrypt images for different reasons : with a chaotic map, a key can be generated in no time, which lead to reasonable computational time. The chaotic map being sensible to tiny variations in initial conditions lead the key sensible to those variations, and consequently, the encryption. It remains, however, a question, which is fundamental in cryptography : is it safe to encrypt images based on chaos theory ? 

In this paper, we will present an algorithm for image encryption, and evaluate it with basic metrics as visual, histogram and correlation analysis. \\
The paper is organized in five sections : the \textit{Section 2} introduces the reader to the algorithm used to encrypt images. A basic analysis of the properties in regards of security is given in \textit{Section 3}. We finally conclude on this paper in \textit{Section 4}. \textit{Section 5} is reserved to mathematical details.


\section{Algorithm}
	\subsection{Encryption}
\footnote{The XOR bitwise operation is denoted as $\oplus$.}
Let's suppose a 2D image $I$ of size $m \times n$, $(m, n) \in \mathbb{N}^2$, with a pixel encoding of $k \in 8\mathbb{N}$ bits per pixel.  \\
In most cases, each pixel of an image is represented as \textit{Red-Green-Blue-alpha (RGBa)} encoded on a byte each (256 possible values for each). 
This lead us to represent the data of the image $I$ as a set $D$ of cardinality $N=m \times n \times k/8$ as follow : 
\begin{equation*}
	D = \{d_1, d_2, ... , d_{N-1}, d_{N}\}
\end{equation*}
with $d_i \in \{0, 1, ..., 255\}, \forall i \leq N$.

\subsubsection{Step 1 : Generation of the encryption key}
The first step is to generate a sequence of numbers $S$ of length $N$ for encryption. This is where we are going to use the chaotic maps, by choosing parameters and initial conditions, we're able to generate large set of real numbers in no time.
We, then, want to define a map $f: U \subset \mathbb{R} \rightarrow \{0,1,...,255\}$, this can be easily done for every chaotic map, an example of this is done \textit{Appendix 2}. \\
As required, we can repeat this step $M \in \mathbb{N}$ times in order to generate an encryption key using the following equation :
\begin{align*}
	K &= f_1(S_1) \oplus f_2(S_2) \oplus ... \oplus f_M(S_M) \\
	&= \{ k_1, k_2, ..., k_{N-1}, k_{N}\}
\end{align*}


\subsubsection{Step 2 : Encryption of the image}
The second step is the encryption of the image $I$. The data is simply encrypted as follow, we denote the data of the encrypted image as a set $E$ of length $N$ :
\vspace{-0.3em}
\begin{align*}
	E &= D \oplus K \\
	&= \{e_1, e_2,..., e_{N-1}, e_N\}
\end{align*}

In fact, we could, if it is required, generate new keys and XOR those keys with the new encrypted image.

\subsection{Decryption}
The main concern while doing cryptography is that we want the encryption to be reversible : indeed, we want to be able to decrypt the encrypted message, image, etc... \\
Mathematically, this means that the map used to encrypt needs to be bijective.
As we proved in the \textit{Appendix 1}, the XOR bitwise operation is a bijective map and the inverse function is the XOR bitwise operation itself. \\
This property of the XOR operation leads us to the following equation to decrypt the encrypted image :
\vspace{-0.4em}
\begin{align*}
	D &= E \oplus K \\
	&= \{d_1, d_2, ..., d_{N-1}, d_N\}
\end{align*}


\section{Results and discussion}
	As we mentionned earlier, images unlike texts have properties which lead traditional algorithm of encryption (mostly block-cypher algorithm like AES-128, RC6...) to some drawbacks. When the algorithm is applied to a pixel at the position $(x,y)$, the values of the four bytes (RGBa) change. These modifications need to appear randomly to the attackers in order to consider the algorithm as a good encryption algorithm. We can analyze the algorithm in several ways, we will present three of them here and we will apply these methods to the proposed algorithm. 

\begin{figure}[!h]	
\centering
\includegraphics[scale=1.15]{tex/img/galois.jpg}
\includegraphics[scale=1.55]{tex/img/grenoble_city.jpg}
\caption{\textit{galois.bmp} and \textit{grenoble\_city.bmp} used for the analysis of the proposed algorithm}
\end{figure}

In order to test the proposed algorithm of the \textit{Section 2}, we are going to encrypt two images : \textbf{galois.bmp} and \textbf{grenoble\_city.bmp}. The first one is a black and white image encoded as an RGBa image, this is interesting because this lead the Red byte, the Green and the Blue one to be equal and the alpha byte set at 255. The second one is a picture of the city of Grenoble in France encoded as an RGBa image. \\

\subsection{Visual analysis}
	Unlike texts, humans can sense patterns in an image by seeing it. The first metrics to measure the efficacity of the algorithm is the visual analysis. \\

As an example of chaotic map, let's consider the well-known discrete logistic map defined as follow :
\begin{equation*}
	r_{n+1} = \mu r_n (1-r_n), \quad \forall n \in \mathbb{N}
\end{equation*}

\begin{figure}[!h]
\centering
\includegraphics[scale=0.8]{tex/img/galois.jpg}
$\rightarrow$
\includegraphics[scale=0.0755]{tex/img/galois1lm.jpg}
$\rightarrow$

\includegraphics[scale=0.0755]{tex/img/galois2lm.jpg} \vspace{1em} \\
\includegraphics[scale=1.01]{tex/img/grenoble_city.jpg}
$\rightarrow$

\includegraphics[scale=0.0955]{tex/img/grenoble_city2lm.jpg}
\caption{Encryption of \textit{galois.bmp} using the logistic map ($\mu=3.89$, $r_0=0.6104$) and then using another XOR logistic map ($\mu=3.79$, $r_0=0.902$) and the encryption of \textit{grenoble\_city.bmp} with the XOR logistic map}
\end{figure}

From the results of the encryption of \textit{galois.bmp}, we can already understand a weakness of the encryption algorithm with the discrete logistic map and black and white images encoded as RGBa images. In fact, the XOR bitwise operation have the following properties, $\forall x \in \mathbb{N}$ :
\vspace{-0.4em}
\begin{align*}
	& x \oplus 0 = x \\
	& x \oplus x = 0
\end{align*}

These properties combined with the high correlation between adjacent pixels lead to some optical illusions which let a human sense the shape of the original image. Obviously, this kind of behaviour need be avoided while doing image encryption. \\
However, this behaviour tends to disappear using multiple discrete logistic maps as it is shown Figure 3. \\

\begin{figure}[!h]
\centering
\includegraphics[scale=1.06]{tex/img/galois.jpg}
$\rightarrow$

\includegraphics[scale=0.1]{tex/img/galoismlm.jpg}
\caption{Encryption of \textit{galois.bmp} using 4 discrete logistic maps with randomly chosen parameters}
\end{figure}

Consequently, using more chaotic maps to generate the encryption key result of more computational time. Inspired by traditional encryption algorithm (AES for example), we could think for bits permutations to avoid the utilisation of multiple chaotic maps. Nonetheless, those permutations required a lot of computational time either, a real benckmarking on various size and type of images need to be done to conclude on the best way to improve the proposed algorithm.

\subsection{Histogram analysis}
	The histogram analysis is a simple analysis which consists to use a bar graph to display the distribution of each byte value. 
A good image encryption algorithm needs to validate two conditions :

\begin{itemize}
\item The graph of the encrypted image needs to be completely different from the one of the original image.
\item The graph of the encrypted image needs to be completely uniform : this means there is the same probability for each byte value to occur in the data.
\end{itemize}

\begin{figure}[!h]
\centering
\includegraphics[scale=0.28]{tex/img/histgrered.jpg}
\includegraphics[scale=0.28]{tex/img/histgregreen.jpg}
\includegraphics[scale=0.28]{tex/img/histgreblue.jpg}
	\caption{Distribution of byte value for each color of the original \textit{grenoble\_city.bmp} image}
\end{figure}

The X-axis of the histogram represents the byte value (between 0 and 255), and the Y-axis is the number of byte in the data for each value.

\begin{figure}[!h]
\centering
\includegraphics[scale=0.28]{tex/img/histgrered_en.jpg}
\includegraphics[scale=0.28]{tex/img/histgregreen_en.jpg}
\includegraphics[scale=0.28]{tex/img/histgreblue_en.jpg}
	\caption{Distribution of byte value for each of the encrypted \textit{grenoble\_city.bmp} with a XOR key of two logistic maps ($\mu_1=3.89$, $r_{01}=0.6104$) and ($\mu_2=3.79$, $r_{02}=0.902$)} 
\end{figure}

In order to study the proposed algorithm, we're going to use the discrete chaotic map defined in the previous section.

We can observe on \textit{Figure 5} that the histograms of the encrypted images are completely different of the histograms of the original image. Furthermore, for each color, the distribution is almost uniform. We can conclude that the proposed algorithm is a good encryption algorithm in regards of the histogram analysis.

\subsection{Correlation between adjacent pixels}
	The main characteristic of images which lead the traditional algorithm of encryption to not be adapted to image encryption is the correlation between adjacent pixels. \\
In order to quantify this phenomenon, we're going to display correlation diagram : random pixels are chosen in the image data, the value of this pixel is displayed on the X-axis, on the Y-axis is displayed the value of the adjacent pixel.

Moreover, the correlation coefficient of 2D variables can be computed following the equation :

\begin{equation*}
	r = \dfrac{cov(x,y)}{\sqrt{D(x)D(y)}}
\end{equation*}

where,
\vspace{-0.5em}
\begin{equation*}
	D(l) = \frac{1}{N} \sum_{i=1}^N \left(l_i - \frac{1}{N} \sum_{k=1}^N l_k \right)^2, \text{ with } l=x \text{ or } l=y
\end{equation*}

and, 
\vspace{-0.5em}
\begin{equation*}
	cov(x,y) = \frac{1}{N} \sum_{i=1}^N \left(x_i - \frac{1}{N} \sum_{k=1}^N x_k \right)  \left(y_i - \frac{1}{N} \sum_{k=1}^N y_k \right)
\end{equation*}


\begin{figure}[!h]
\centering
\includegraphics[scale=0.5]{tex/img/corregal.jpg}
\caption{Correlation diagram ($r = 0.981564$) for 1000 randomly chosen points of the original image \textit{galois.bmp}}
\end{figure}

To be considered as a good image encryption algorithm, the correlation diagram should not contain discernable patters and the correlation coefficient should be close enough to 0. 
\begin{figure}[!h]
\centering
\includegraphics[scale=0.5]{tex/img/correegal.jpg}
\caption{Correlation diagram ($r = 0.083521$) for 1000 randomly chosen points of the encrypted image \textit{galois.bmp} with a XOR key of two logistic maps ($\mu_1=3.89$, $r_{01}=0.6104$) and ($\mu_2=3.79$, $r_{02}=0.902$)} 
\end{figure}

We can observe on \textit{Figure 7} that the correlation diagram of the encrypted image, with the proposed algorithm and the discrete logistic map, validate the requirement cited above. The proposed algorithm can be considered as a good encryption algorithm regarding the correlation of adjacent pixels analysis. 

\section{Conclusion}
	In this paper, we have shown that chaos-based cryptography can be suitable to image encryption.
Indeed, we have proposed a naive implementation which lead to great results in various analysis tests such as visual, intensity and correlation between adjacent pixels. 
However, as we shown earlier in the paper, multiple chaotic maps are needed to get safe encryption, this could lead to larger computational time. \\

The proposed algorithm could be a basis to more complex algorithms for image encryption, for example, to reduce compute time, the representaion of floating-point number as 32 bits (or double floating-point number as 64 bits) can be use to avoid the encryption byte per byte as it is done in this paper.

\section{Appendix}
\subsection{Demonstration of the bijectivity of the bitwise operation XOR}
	In order to prove the bijectivity of the bitwise operation XOR, we need to prove the injectivity and the surjectivity of the operation. The surjectivity is pretty trivial, however, to prove that the XOR operation is a one-to-one mapping seems a little bit tricky. 

Let's denote the XOR operation as a function $f:\mathbb{N} \times \mathbb{N} \rightarrow \mathbb{N}$. We want to prove that $f(x,y)=f(x',y) \Rightarrow x=x'$ (we will suppose $y$ being a constant, \textit{the encryption key}). Let's take $(x,y) \in \mathbb{N}^2$ as follow :

\begin{equation}
	x = \sum_{i=0}^{\infty}\alpha_i 2^i \qquad \text{where } \forall i \in \mathbb{N}, \alpha_i \in \{0,1\}
\end{equation} 

and,

\begin{equation}
	y = \sum_{i=0}^{\infty}\beta_i 2^i \qquad \text{where } \forall i \in \mathbb{N}, \beta_i \in \{0,1\}
\end{equation} 

Analytically, we can rewrite the $f(x,y)$ function as :

\begin{equation}
	f(x,y)=\sum_{i=0}^{\infty} |\alpha_i - \beta_i| 2^i
\end{equation}

Hence, if we take $x'$ as :

\begin{equation}
	x' = \sum_{i=0}^{\infty}\gamma_i 2^i \qquad \text{where } \forall i \in \mathbb{N}, \gamma_i \in \{0,1\}
\end{equation} 

The injectivity condition leads us to :

\begin{equation}
	f(x,y)=f(x', y)=\sum_{i=0}^{\infty} |\alpha_i - \beta_i| 2^i=\sum_{i=0}^{\infty} |\gamma_i - \beta_i|2^i
\end{equation}

Hence, by polynomial identification, we got :

\begin{equation}
	|\alpha_i - \beta_i| = |\gamma_i - \beta_i|
\end{equation}

This condition implies $x=x'$ which ends the demonstration.
Note that from the analytical definition of the XOR operation, it is clear that the inverse of the XOR operation is the XOR operation itself.

\subsection{Mapping to a byte of the logistic map}
	For this paragraph, we will denote the floor function of $x$ as $\left\lfloor x \right\rfloor$. \\

Let's consider the discrete logistic map defined as : 

\begin{equation*}
	r_{n+1} = \mu r_n (1-r_n), \quad \forall n \in \mathbb{N}
\end{equation*}

Since at any point, we have $0 \leq r_n \leq 1$, a map $f:U \subset \mathbb{R} \rightarrow \{0,1,...,255\}$ which satisfy the requirement for the paper, can be the following map :

\begin{equation*}
	\forall x \in U, f(x) = \left \lfloor x * 255 \right \rfloor
\end{equation*}

\end{document}

